\chapter{元素从哪里来，又去了哪里}

\begin{kaichang}
想象两样东西。

第一样，是一枚金戒指——也许就在你家里某个长辈的手指上。金子不会生锈，不会腐烂，放上几百年还是黄澄澄的。可是，地球上的金子是从哪里来的？金矿挖完以后，我们能像种庄稼一样“再种一批”吗？能像在工厂里生产塑料那样“合成”一批吗？

第二样，是一张被你烧掉的纸（请只想象，别动手）。火苗蹿过，纸变成了灰烬和一股烟，轻飘飘的。纸到哪里去了？有人说“烧没了”。真的没了吗？

一个“从哪里来”，一个“去了哪里”。这两个问题听起来平常，答案却要动用全书最宏大的尺度：从宇宙诞生后的前三分钟，到恒星的心脏，再到你脚下的泥土和你碗里的米饭。这一章走完，你会知道那枚金戒指的来历，远比你想象的浪漫。
\end{kaichang}

\section{先盘点：我们身边的元素库存}

在回答“从哪里来”之前，先看看“现在都在哪”。我们可以给几个熟悉的地方做一份“元素库存盘点”——按质量百分比，排出前几名：

\begin{table}[htbp]
\centering
\begin{tabular}{lllll}
\toprule
排名 & 地壳 & 海水（溶解部分之外含水） & 大气 & 人体 \\
\midrule
第 1 名 & 氧（约 46\%） & 氧（约 86\%） & 氮（约 78\%） & 氧（约 65\%） \\
第 2 名 & 硅（约 28\%） & 氢（约 11\%） & 氧（约 21\%） & 碳（约 18\%） \\
第 3 名 & 铝（约 8\%） & 氯、钠等 & 氩（约 0.9\%） & 氢（约 10\%） \\
第 4 名 & 铁（约 5\%） & 其余几十种 & 二氧化碳等 & 氮（约 3\%） \\
\bottomrule
\end{tabular}
\end{table}

这张表里藏着好几个意外。地壳里最多的不是铁也不是碳，而是氧——沙子、岩石、黏土的主要成分都是硅和氧搭起来的骨架。人体按质量算第一名居然也是氧（因为人体约六到七成是水），而碳这个“生命骨架元素”只排第二。大气则是氮气的天下，氧反而只有五分之一。

还要注意：这些是平均数。你的手机里有几十种元素，包括地壳里只占百万分之几的镓、铟、钽；深海热泉边的细菌体内，钼和钨的含量比陆地生物高得多。元素的分布极不均匀——有的地方富集成矿，有的地方稀薄到要用一吨矿石换几克金属。这份“不均匀”，后面会变成一个大问题。

\section{化学反应动不了原子核}

第 7 章讲过，元素的身份证是原子核里的质子数：8 个质子就是氧，79 个质子就是金，多一个少一个都不行。第 9 到 11 章又讲了，化学反应的全部“剧情”——成键、断键、得失电子——都发生在核外的电子身上，原子核自始至终坐在后台，纹丝不动。

把这两件事拼起来，就得到本章第一条规则：

\begin{qiaoliang}
\textbf{在一切普通化学反应中，元素的种类和原子总数都不会改变。}反应前有多少个碳原子，反应后还是多少个碳原子；它们只是换了邻居、换了组合方式。想让一种元素真正变成另一种元素，必须改写原子核本身——那是核反应，需要的能量比化学反应高出几十万到几百万倍，绝不是点一把火、倒一瓶药水能做到的。
\end{qiaoliang}

这就是“烧掉的纸去了哪里”的一半答案：纸里的碳、氢、氧原子一个也没消失。它们大部分和氧气结伴，变成二氧化碳和水蒸气散进了空气，小部分留在灰烬里。你之所以觉得“没了”，是因为产物从看得见摸得着的固体，变成了看不见的气体——这是形态的消失，不是原子的消失。

生活里处处可以验证这条规则，只要你把“藏起来的部分”也计入账内。铁钉生锈后变重了，多出来的质量不是凭空长的，是从空气里“抓”来的氧；银首饰放久了发黑，是银和空气里微量的硫化氢结伴成了硫化银，银原子一个没少，只是换了搭档；你吃进去的米饭，一部分碳变成你呼出的二氧化碳，另一部分砌进了你的骨骼和肌肉——\textbf{没有任何一个原子是在你体内“生产”出来的}，你只是一座繁忙的原子中转站。这也顺便解释了为什么人必须从食物里摄取铁、锌、碘这些“矿物质”：身体能组装分子，却变不出元素，缺了就只能从外界补给现成的原子。

下面这个实验，可以让你亲手“抓住”这些不肯消失的原子。

\begin{shiyan}
\textbf{密封袋里的守恒}

材料：一个能密封的透明塑料袋（带密封条的保鲜袋最好）、一小勺小苏打（碳酸氢钠）、一小杯白醋、一张厨房纸巾、一台厨房电子秤（能读到 1 克即可）。

步骤：
\begin{enumerate}[itemsep=1pt]
\item 把小苏打倒在纸巾上，包成小纸包，放进塑料袋。
\item 把白醋倒进袋子里，注意别让醋立刻碰到纸包。
\item 排出大部分空气后封紧密封条，把整个袋子放在电子秤上，记下质量。
\item 不动秤，隔着袋子捏破纸包、晃一晃，让小苏打和醋反应。观察：袋子慢慢鼓起来，里面满是气泡。
\item 反应结束、袋子鼓起后，再读一次质量。
\end{enumerate}

观察点：反应前后电子秤的读数几乎不变（差别通常小于秤本身的误差）。袋子里明明发生了激烈的反应，生成了大量气体，可总质量没有变——因为生成的二氧化碳分子全被关在了袋子里。如果把袋子打开再称，质量会变小，不是原子消失了，而是气体“越狱”了。

安全提示：这个实验很安全，但请在水槽边操作，免得袋子漏气喷出醋水；不要用玻璃瓶代替塑料袋密封反应，气体膨胀可能撑碎玻璃。
\end{shiyan}

\section{那原子核里的元素，是谁改写的？}

化学反应造不出金，也毁不掉碳。那么问题原封不动地回来了：今天世界上这一百多种元素，最初是从哪儿来的？答案是：它们分三批出生，出生在三种完全不同的“工厂”里。

\textbf{第一家工厂：大爆炸。}第 8 章提过，宇宙诞生于约 138 亿年前。刚诞生的宇宙热得无法形容，随着膨胀冷却，质子和中子才来得及“冻结”出来。在最初大约三分钟里，一小部分质子和中子碰巧撞在一起、粘成了氦核，还有极少量粘成了锂。然后宇宙就冷得、稀得再也聚不动了。所以第一批元素简单得可怜：按质量算约四分之三是氢，四分之一是氦，再加一丁点锂。你身体里的每个氢原子，几乎都是那三分钟里出生的，年纪 138 亿岁。

\textbf{第二家工厂：恒星。}大爆炸造不出碳、氧、铁——没有它们就没有行星和生命。这批元素要等几亿年后，引力把氢氦气体云拧成恒星之后才登场。恒星靠核聚变发光：核心上千万度的高温高压下，氢核聚成氦核，聚变放出的能量顶住引力的挤压，让恒星稳定燃烧几十亿年。氢烧完后，核心收缩、变得更热，氦又开始聚成碳和氧；更大的恒星还能一层层接着烧：碳烧成氖、镁，再烧成硅，最后烧成铁。像洋葱一样，恒星晚年内部一层一种“燃料”。

这套流水线有个反直觉的排班表：越往后的“燃料”，烧得越快。太阳的氢够烧约 100 亿年；大质量恒星烧氢也要几百万到上千万年，可氦阶段只剩几十万年，碳阶段几百年，硅阶段——只有大约一天。仿佛一座工厂前面几道工序慢悠悠，最后一道工序一天之内冲刺完工，然后整座工厂轰然停工。所以宇宙里轻元素堆积如山，而铁以上的元素永远是零头：氢和氦占了可见物质的 98\% 以上，其余一百多种元素加起来不到 2\%。

但这条路到铁就断了。原因在第 8 章见过一面，这里说透：

\begin{qiaoliang}
原子核里，质子互相排斥，靠强核力把它们粘在一起。把分散的核子捏成核时会放出能量，这就是“结合能”。平均到每个核子头上，结合能最高的元素恰好是铁附近——铁核是“最紧实、最省心”的状态。所以：比铁轻的元素聚变，总能量降低，放能；比铁重的元素聚变，反而要倒贴能量。铁是恒星核燃烧的终点站：恒星核心一旦烧成铁，就再挤不出能量来对抗引力，大质量恒星随即坍缩、爆发。
\end{qiaoliang}

\textbf{第三家工厂：恒星的死亡与相撞。}比铁更重的元素——银、金、铂、铀、碘——只能靠更暴烈的场面来造。一种是大质量恒星死亡时的超新星爆发：坍缩瞬间释放的巨大能量把铁核打碎又重捏，硬把核子塞进更重的核里。另一种是两颗中子星（恒星死亡留下的、一茶匙就有几亿吨重的致密残骸）相互绕转、最终撞在一起：碰撞把大量中子泼洒到周围的原子核上，核一口气吞下几十个中子，再经过一连串衰变，稳定成金、铂这样的重元素。你那枚金戒指里的每个金原子，都来自某次这样的宇宙车祸，时间在太阳系诞生之前。

\section{用符号把三种工厂写下来}

现在轮到符号登场。核反应式和化学反应式长得很像，规矩却不同：化学方程式两边原子种类守恒；核反应式两边守恒的是质子总数和质量数总数，元素种类恰恰是要变的。看这三条：

氢聚变（太阳的主要能源，简化写法）：
\[4\,\chem{^{1}_{1}H} \rightarrow \chem{^{4}_{2}He} + \text{能量}\]

碳的诞生（红巨星里的“三氦反应”）：
\[3\,\chem{^{4}_{2}He} \rightarrow \chem{^{12}_{6}C} + \text{能量}\]

金的一种出生路径（中子星并合中，一个核子核吞下中子后衰变）：
\[\chem{^{197}_{80}Hg} \rightarrow \chem{^{197}_{79}Au} + \chem{^{0}_{+1}e}\]

注意第三条：汞（80 号）衰变后真的变成了金（79 号）——中世纪炼金术士的梦想，在物理上并非不可能。那为什么不造金子发财？读完本章的误解框你就明白了。

顺带算一笔账，体会一下“每个原子都大有来头”。一枚 \ud{5}{g} 的金戒指，含多少金原子？金的摩尔质量约 \ud{197}{g\cdot mol^{-1}}，所以 $5/197 \approx 0.025\ \mathrm{mol}$，乘上阿伏伽德罗常数：
\[0.025 \times 6.02\times 10^{23} \approx 1.5\times 10^{22}\ \text{个金原子}\]
一万五千亿亿个原子，每一个都是在几十亿年前某次中子星相撞的烈焰里锻造出来，又搭乘气体云、陨石、岩浆，辗转几十亿年才躺进这枚戒指。浪漫吗？这是算出来的，不是诗人的修辞。

\section{科学家凭什么知道恒星在造元素？}

我们没钻进过任何一颗恒星的核心，凭什么一口咬定元素是在那里造出来的？因为有两类铁证：一类来自光，一类来自“看得见的预言”。

\begin{zhengju}
\textbf{第一条线索：先在太阳上发现，后才在地球上找到的元素。}1868 年，日全食期间，天文学家把分光镜对准太阳边缘的日珥，看到一条明亮的黄色谱线——它和当时地球上所有已知元素的谱线都对不上。有人大胆宣布：太阳里有一种地球上还没发现的元素，就用希腊语“太阳”给它起名：氦（helium）。这个宣布相当冒险：万一是仪器误差或已知元素的某种未知状态呢？争议持续了近三十年。直到 1895 年，化学家拉姆齐在地球的一种铀矿加热后放出的气体里，找到了同一条黄色谱线——氦真的存在，而且地球上有。这件事确立了分光术的地位：\textbf{每种元素的发光谱线独一无二，像条形码}，光就是远方物质的成分报告。从此，天文学家读出：太阳约四分之三是氢、四分之一是氦——和大爆炸的“第一批产品”配方吻合；老年恒星的光谱里碳、氧、钠、铁的谱线随年龄增强，和“恒星边烧边造”的剧本吻合。

\textbf{第二条线索：一次被全球围观的“预言兑现”。}1957 年，伯比奇夫妇、福勒和霍伊尔发表论文，系统提出重元素在恒星与超新星中合成的理论，其中预言：金、铂这类最重的元素需要在“中子暴雨”中快速捕获中子才能形成，而中子星相撞应该是最理想的场景之一。预言之后是长达六十年的等待。2017 年 8 月，引力波探测器捕捉到约 1.3 亿光年外两颗中子星相撞的涟漪；几乎同时，全球几十台望远镜转向同一方向，在可见光和红外波段看到一团快速膨胀变暗的余辉——“千新星”。余辉的光谱随时间变化的方式，和“大量重元素正在形成并衰变”的计算结果对得上。据估算，这一次碰撞抛出的金就有好几百个地球质量。从光谱条形码到引力波，两条完全独立的证据链指向同一结论：恒星和恒星的死亡，就是元素的铁匠铺。
\end{zhengju}

\section{元素不出门，只在地球上搬家}

元素不会凭空增减，但绝不会躺着不动。它们在岩石、海洋、大气和生物之间不停地流转，科学家把这些流转路线叫作“循环”。挑三条对生命最重要的看一眼：

\begin{itemize}[itemsep=2pt]
\item \textbf{碳循环}：二氧化碳被植物光合作用固定进糖和木质素，动物吃植物，动植物的呼吸、死亡与腐烂又把碳还给大气和土壤；一部分碳沉到海底变成石灰岩，要等地壳运动把它抬回地面。这条循环里，碳原子的“通勤时间”从几天（呼吸）到上亿年（石灰岩）不等。近两百年来，人类把煤和石油里封存了几亿年的碳快速烧回大气，相当于把慢车道上的货一次性全搬到快车道——碳原子还是那个碳原子，但它们在大气里的浓度变了，气候随之改变。第 54 章讲光合作用时会回到这条循环的“入口”。
\item \textbf{氮循环}：空气里五分之四是氮气，可氮气分子太稳定，绝大多数生物用不上。只有少数细菌和闪电能把氮“固定”成可用的含氮化合物，植物吸收，动物吃植物，微生物再把动植物的遗体和排泄物分解，一部分氮变回氮气回到天空。一百多年前，人类发明了工业固氮（哈伯法），用高温高压把氮气和氢气压成氨——今天全球约一半人口吃的粮食，靠的是固氮工厂养活的庄稼。这是化学改变文明的最硬核案例之一，第 14 章细讲。
\item \textbf{磷循环}：磷没有气态形式，只在岩石、土壤、水体和生物间慢慢搬运：岩石风化释放磷酸盐，植物吸收，沿食物链传递，最后随遗体沉入水底，要等千万年级别的地质过程才能回来。磷几乎是“单程车票”，所以农田要施磷肥，而生活污水把大量磷冲进湖泊，又会催生藻类疯长、耗尽氧气（富营养化）。
\end{itemize}

请注意这三条循环共同的画风：\textbf{没有一条循环涉及元素的创生或消灭}。所谓的“循环”，全是一群亘古不变的原子在换搭档、换住所。驱动它们搬家的能量，最终几乎都来自太阳（第 54 章往后的代谢篇章，就是讲生命怎样借助这股能量，把原子的搬家组织成“活着”这件事）。

\section{既然不会消失，为什么还会“短缺”？}

这时你可能想抬杠：既然元素不会消失，新闻里为什么总说锂、钴、稀土“紧缺”？问得好，这恰好是本章在工业上的应用点。

短缺从来不是“原子没了”，而是两个问题。第一是\textbf{浓度}：一部手机的锂电池里大约只有几克锂，一辆电动车约几公斤；这些锂用完后并没有消失，而是分散在几亿块废旧电池里，和塑料、铝、电解液混在一起。把分散的元素重新聚拢、提纯，要花能量和钱——稀释容易浓缩难，这正是第 27 章要讲的热力学第二定律在资源问题上的现身。第二是\textbf{分布}：地壳里的稀土其实不算稀，但富集成值得开采的矿床的地方很少；钴的储量一多半集中在刚果（金），锂的高浓度盐湖集中在南美“锂三角”和我国青藏高原。元素全球化地循环，矿床却极度挑剔地分布，于是供应链成了地缘政治问题。

解决办法也写在本章的逻辑里：既然原子不灭，\textbf{回收就是“城市矿山”}。一吨废旧手机线路板里的金含量，可以是一吨金矿石的几十倍。日本曾为东京奥运会奖牌回收了约 8 万吨废旧小家电，提炼出约 \ud{32}{kg} 金、\ud{3500}{kg} 银、\ud{2200}{kg} 铜。让元素在经济系统里多转几圈，比去地壳深处追新矿，往往更省能量、更少污染。

\section{这条规则在哪里会失灵}

诚实地说一下边界。“元素不生不灭”只在\textbf{化学和生物过程}的尺度上成立，超出范围就要打补丁：

\begin{itemize}[itemsep=2pt]
\item \textbf{核反应不守这条}。第 8 章的放射性衰变会把铀慢慢变成铅，医院的核医学每天都在把一种元素变成另一种；太阳此刻正把氢变成氦。只是这些过程的能量尺度，与普通生活和普通化学隔着好几个数量级。
\item \textbf{“库存表”不是永恒的}。地壳、大气的元素比例在地质时间里会变：二十多亿年前大气里几乎没有氧，是蓝细菌的光合作用用几亿年把氧“造”出来的（氧原子不新，是它们从水里释放出来这件事新）。放射性元素如铀-238 则在不可逆地减少，半衰期约 45 亿年——太阳系的铀已经烧掉了将近一半。
\item \textbf{极轻微的“质量—元素”换算}。任何放能反应都会让产物总质量轻一点点（见挑战框），只是化学反应的这点质量差小到天平永远称不出来，所以“原子守恒”在化学里是极好的近似。说“近似”，不是说它错，而是知道它的说明书上写着适用范围。
\end{itemize}

\begin{wuqu}
\textbf{“既然核反应能变元素，那造黄金发财不就行了？”}

想法没错，账算不过来。实验室里确实用粒子加速器把铋或汞变成过金：1980 年，美国劳伦斯伯克利实验室的团队用加速器轰击铋靶，真的得到了微量金。代价呢？耗费的电费和设备折旧折算下来，每造出 1 克金的成本是金价本身的几百万倍，而且产物常常混着放射性同位素，得等它衰变安静了才敢碰。反过来，历史上炼金术士用炉火、药水折腾了几百年，一原子金也没造出来——因为他们用的全是化学手段，而化学手段碰不到原子核。这个误解值得记住，因为它把两件事分得很清：\textbf{“物理上可能”和“经济上可行”之间，隔着能量、成本和规模的万丈深渊}。元素守恒被核物理打破了一角，但被经济规律补上了十堵墙。
\end{wuqu}

\section{回到金戒指和那张纸}

现在兑现开场的两个承诺。

金戒指：它的每个原子出生于太阳系形成之前某次中子星相撞（或超新星爆发），随着抛撒出的碎屑混入气体云；约 46 亿年前，这朵云收缩成太阳系，金原子跟着重元素一起沉进正在成型的地球；地球早期还是熔融状态时，大部分金沉向了地心，只有一小撮留在地壳，被后来的热液富集成矿；矿工把它挖出，工匠把它打成戒指。它不会被“种”出来，也不会在化学反应中增减半分——地球上的金，基本上就是那批，用一点，散一点，只能靠回收聚拢。

那张烧掉的纸：碳和氢多数变成了 \chem{CO_2} 和 \chem{H_2O} 进了空气。接下来的剧情你可以自己续：某个二氧化碳分子可能下周被一片叶子抓住，变成葡萄糖；可能溶进海水，被浮游生物做成碳酸钙外壳，沉到海底，几百万年后成为石灰岩，被炸开、烧成水泥，砌进某栋楼的墙里。原子不死，只是永远在赶下一场。

\section{伏笔：原子不变，搭档常新}

这一章我们盯着原子核，得到的是一副近乎静态的画面：元素守恒，原子永恒。可生活里的物质明明瞬息万变——纸会燃烧，铁会生锈，食物会变成你的一部分。变化的全部秘密，不在原子的生死，而在原子之间\textbf{连接方式}的重组：谁和谁结伴、以多大的力气拉扯、重组时能量账怎么算。从下一章开始，镜头从原子核转向原子之间——第 17 章，我们去拆“化学键”这个被误会最深的概念。

\begin{tiaozhan}
\textbf{恒星哪来的那么多能量？一笔质量账。}聚变的能量来自质量本身：4 个质子聚成 1 个氦核后，产物的质量比反应前少了约 $0.7\%$。按 $E=mc^2$ 换算，1 克氢聚变损失约 \ud{0.007}{g} 质量，释放能量约 $0.007\times 10^{-3}\ \mathrm{kg}\times (3\times 10^{8}\ \mathrm{m/s})^2 \approx 6\times 10^{11}\ \mathrm{J}$——相当于燃烧约 20 吨煤。太阳每秒把约 6 亿吨氢烧成氦，质量每秒“亏掉”约 400 多万吨，全靠光速的平方这个天文数字撑着，才烧了 46 亿年还没烧完（才烧了总氢量的百分之几）。也正因为 $c^2$ 如此巨大，化学反应里的质量变化（几十亿分之一量级）才完全称不出来。跳过这段不影响主线，但它能帮你把第 8 章的结合能、本章的“铁是终点”和下一篇的能量账串成一条线。
\end{tiaozhan}

\section*{练一练}

\begin{sikaoti}
\item 画一幅“元素三工厂”流程图：用三个方框（大爆炸、恒星聚变、恒星死亡/相撞）和箭头，标出氢、氦、碳、氧、铁、金分别从哪个工厂出厂。图旁注明：方框的大小不代表真实时间和空间尺度。
\item 一根木柴燃烧后，灰烬的质量远小于木柴。有同学说“这说明质量守恒在燃烧中不成立”。用粒子层面的语言反驳他，并设计一个（只用厨房材料的）改进方案来验证你的反驳。
\item 查一下你呼吸一次大约交换 \ud{500}{mL} 空气。结合大气约 78\% 是氮气，估算你一天吸入的氮气大约有多少摩尔（常温常压下 1 摩尔气体约 \ud{24}{L}）。这些氮气大部分“原样呼出”，这说明了氮循环的什么特点？
\item 参考正文中金戒指的算法，估算你体内约 \ud{4}{g} 铁（主要在血红蛋白里）包含多少个铁原子，并说明这些原子最可能出自哪一家“工厂”（提示：铁在元素合成中的特殊地位）。
\item 某广告称某保健品“含有人体能自己合成不了、必须额外补充的特殊活性元素”。用本章内容指出这句话里合理的部分和偷换概念的部分（人体不能合成的是元素本身，还是某些含这些元素的分子？）。
\item 画一条“一个碳原子的万年旅行”路线图：从大气出发，至少经过植物、动物、土壤或海洋中的三站，再回到大气。在每站标注这个原子的“搭档”变成了什么（\chem{CO_2}、葡萄糖、碳酸盐……）。
\end{sikaoti}
